\paperabstract{%
在大国科技竞争和出口管制持续强化的背景下，半导体相关产品进口来源结构已成为评估产业链供应链安全的重要切入点。本文基于 CEPII BACI HS07 数据库 \citep{cepii_baci}，构建 2008--2024 年中国自全球进口四类半导体相关 HS6 产品的出口国--产品--年份平衡面板数据，从描述统计、风险指数和统计检验三个层次考察政策冲击下的进口来源结构演化。研究首先从进口规模、美国份额、主要来源国和 HHI 集中度四个角度刻画结构变化事实；其次构建半导体进口供应链风险指数 SIRI，将来源集中、政策暴露、替代不足和结构波动纳入统一评分框架，按产品--年份计算 0--100 综合风险分值；最后结合 OLS 固定效应回归（进口额）和 PPML 估计（进口份额）以及多项稳健性检验，讨论政策节点下的统计证据边界。结果显示，2017--2024 年间，半导体制造设备（848620）的美国进口份额从 27.92\% 降至 9.88\%，处理器及控制器（854231）和存储器（854232）同样下降，但其他电子集成电路（854239）方向相反。2024 年 SIRI 排序显示，854239 和 854232 的综合风险相对更高，主要驱动因素是来源集中和替代不足而非政策暴露——提示高对美依赖不等于高综合风险。统计检验揭示了一个层次化的管制传导过程：2018 年实体清单扩容后美国来源绝对进口额显著上升（$p=0.002$），但份额未显著下降，反映了企业的"恐慌性囤货"行为；2022--2023 年出口管制新规成为转折点，PPML 估计显示美国来源份额下降约 25--27\%（$p<0.001$），绝对进口额于 2023 年开始显著相对下滑。产品层面异质性显著：处理器类产品 2018 年后进口额先升后降，存储器类产品 2018 年后即大幅下降，其他集成电路类产品则持续上升。进一步地，本文在未接入 GDELT 的 BACI-only 设定下，将产品池扩展至 20 个半导体相关 HS6 编码，构造年度进口来源网络，并进行下一年 SIRI 风险预测的原型验证。本文为关键半导体产品进口安全监测提供了可解释、可复现的统计建模框架，并讨论了结论的适用范围与局限。
}{大国博弈；半导体；进口来源结构；供应链风险；统计建模}
